\section{Evaluation}\label{sec:experiments}

The six public DAGs are the released instances of Problem A, ``Intra-kernel
Scheduling on a General-Purpose Neural Network Processor,'' from the 2025
Huawei Cup China Graduate Mathematical Contest in Modeling. They cover
convolution, FlashAttention, and matrix multiplication and range from 1,716 to
36,086 nodes. All experiments use the capacities and
evaluator of Section~\ref{sec:formulation}; every production artifact has zero
violations. Official artifacts are used only for final comparison, and no order
generator reads an official schedule. The scalable solver is the sole production
default; repair and exact planning are reported separately.

\subsection{Public P2 and P3 Results}\label{sec:exp-p2}

Table~\ref{tab:public} and Figure~\ref{fig:public-results} give the primary
official-artifact comparison. P2 has five strict wins and one tie (median
reduction $0.866\%$, maximum $9.08\%$); P3 time has five wins and one loss
(median improvement $3.77\%$), with Conv\_1 the sole $4.23\%$ regression.
Complete P1/P2/P3 metrics, counts, tie-break times, and validation states are in
Supplementary Tables~S5--S7.

\begin{table}[t]\centering\small
\caption{Canonical production solver versus official P2/P3 artifacts.}
\label{tab:public}
\begin{tabular}{lrrr}\toprule
Instance & P2 ours & P2 official & P3 time: ours/official\\\midrule
Conv\_0 & \textbf{66,828} & 73,500 & \textbf{515,634}/535,312\\
Conv\_1 & \textbf{72,734} & 73,240 & 1,118,687/\textbf{1,073,322}\\
FA\_0 & \textbf{3,584} & 3,692 & \textbf{36,344}/46,761\\
FA\_1 & \textbf{32,512} & 32,840 & \textbf{152,364}/193,059\\
Matmul\_0 & \textbf{34,688} & 34,944 & \textbf{186,820}/194,331\\
Matmul\_1 & 460,800 & 460,800 & \textbf{1,771,383}/1,800,218\\\bottomrule
\end{tabular}\end{table}

\begin{figure*}[t]
\centering
\includegraphics[width=0.9\textwidth]{\ksFigure{headline_reductions.pdf}}
\caption{Production results on all six public instances: P2 traffic reduction
against official artifacts and P3 execution-time reduction. Repair and the
fixed-order oracle are excluded.}
\Description{Two horizontal panels showing five P2 wins and one tie and five
P3 wins with one regression.}
\label{fig:public-results}
\end{figure*}

A controlled $H=0$ ablation is descriptively nonadditive but is not a complete
factorial design. Conv\_1 remains 214 bytes above the actual predecessor, and
both Matmul cases are unchanged. We therefore interpret complete configurations,
not isolated main effects or synergy; Supplementary Table~S8 and Figure~S1 give
the full path.

\begin{figure*}[!t]
\centering
\includegraphics[width=0.92\textwidth]{\ksFigure{order_headroom.pdf}}
\caption{Logical L1 residency on Conv\_Case0 under two certified topological
orders, decomposed into backed and generated live bytes. The dependency-frontier
order has a larger geometric overflow area but a $29.6\%$ lower certified
fixed-order traffic optimum (57,408 versus 81,504 bytes). Geometric pressure is
therefore not a reliable order-ranking surrogate under asymmetric cost.}
\Description{Two aligned residency traces compare the legacy P1 and
dependency-frontier orders against the 4-KiB L1 capacity.}
\label{fig:order-headroom}
\end{figure*}

\subsection{\texorpdfstring{$\mathrm{Tr}=\mathrm{Vol}+\mathrm{Dt}$}{Tr=Vol+Dt} Accounting}
\label{sec:exp-decomposition}

The cases span generated-only, mixed, and backed-only regimes
(Figure~\ref{fig:vd-plane}). Every strict P2 win reduces $\mathrm{Vol}$;
Conv\_1 wins despite increasing $\mathrm{Dt}$ by 171 bytes, while Matmul\_0 wins
with $\mathrm{Dt}=0$. No single composition story explains the results.
Per-case accounting is in Supplementary Table~S9 and Figure~S2.

\begin{figure}[t]
\centering
\includegraphics[width=\columnwidth]{\ksFigure{vd_plane.pdf}}
\caption{Public cases on the normalized $(\mathrm{Vol},\mathrm{Dt})$ plane.
Arrows point from official to scalable artifacts; dashed diagonals are
iso-traffic lines. The cases span generated-only, mixed, and backed-only regimes.}
\Description{Per-case points and arrows across diagonal iso-traffic lines.}
\label{fig:vd-plane}
\end{figure}

Two nonuniform repair studies hold the planner fixed and change only a legal
order, reducing Conv\_0/Conv\_1 by 1,296/1,794 bytes. Post-hoc inspection
associates accepted moves with advancing backed \textsc{alloc}/\textsc{copy-in}
chains, but $\mathrm{Vol}$ and $\mathrm{Dt}$ change together, so this does not
isolate a cause. Heterogeneous budgets, negative FlashAttention probes, and
untested Matmul appear in Supplementary Table~S11.

\subsection{Exact-to-Heuristic Evidence}\label{sec:exp-exact}

The Conv\_0/frontier fixed-order certificate is 57,408 bytes, $14.1\%$ below
the scalable plan under that order. The same oracle certifies 81,504 bytes under
the legacy P1 order, isolating a $29.6\%$ order difference. FA\_0 supplies the
third certificate and covers a second public instance. Other cases have positive
gaps, timeout, or were not run (Table~\ref{tab:exact}) and are not traffic-optimal.
Certificates do not optimize P2 spill-count or time tie-breakers. The full chain,
failures, and machine-artifact scope are in Supplementary Secs.~3.3--3.5 and
Table~S12.

\begin{table}[t]\centering\scriptsize
\caption{Repair case studies and fixed-order exact evidence. A dash is no
completed result, not zero traffic.}
\label{tab:exact}
\begin{tabularx}{\columnwidth}{lrrrX}\toprule
Instance & Scalable & Repair & Fixed-order & Status\\\midrule
Conv\_0 & 66,828 & 65,532 & \textbf{57,408} & Traffic certificate\\
Conv\_1 & 72,734 & 70,940 & -- & Packing timeout\\
FA\_0 & 3,584 & probe & \textbf{3,584} & Traffic certificate\\
FA\_1 & 32,512 & probe & 32,512 & Feasible; LB 31,936\\
Matmul\_0 & 34,688 & -- & 34,816 & Feasible; LB 29,952\\
Matmul\_1 & 460,800 & -- & -- & Not run\\\bottomrule
\end{tabularx}\end{table}

\paragraph{Order headroom versus pressure.}
Conv\_0/frontier has larger logical L1 overflow area than the legacy P1 order
but $29.6\%$ lower certified traffic. Live/overflow-area surrogates therefore
cannot reliably rank orders under asymmetric cost (Figure~\ref{fig:order-headroom}).

\subsection{Robustness, Negative Results, and Limitations}\label{sec:limitations}

\paragraph{Capacity sweep.}
At every scanned feasible capacity point in the Conv\_0/L1/fixed-$H=0$ sweep,
the selected configuration never loses to the cp-free-first order comparator;
where both traffic values are nonzero, the ratio grows from $1.35\times$ to
$124\times$. This is a controlled single-instance, single-cache boundary, not
monotonicity of the absolute difference or cross-operator robustness. The full
nine-point data, infeasible boundary, and curve are in Supplementary Table~S13
and Figure~S4.

\paragraph{Internal synthetic boundary.}
On 36 standard re-evaluations, production ties its predecessor candidate subset;
cp-free-first is also identical to included id-raw on 36/36 instances, so both
zero-loss results are principally portfolio-monotonicity checks. We observe zero
losses against four nonincluded fixed implementations; all eight 17-node oracle
instances attain the fixed-order traffic optimum under their selected order.
This supports non-regression and small-scale same-order agreement, not new
generalization over the predecessor. The 36+8 per-instance rows and 252
run-time validations are reported in Supplementary Tables~S14--S17; artifact
locations and reproduction commands are maintained in the repository guide.

\paragraph{Scope threats.}
The public evaluation contains only six DAGs and the exact study is incomplete.
The benchmark uses static \textsc{copy-in} labels, requires residency at
\textsc{free}, and exposes no read/write roles, precluding claims about dynamic
dirty-state evolution. P1 is a logical footprint, not physical P2 residency;
CP-SAT may time out, and scalable victim policies remain heuristic. All claims
are limited to the stated evaluator and artifact contract.
